\subsection*{Factores limitantes}
\addcontentsline{toc}{subsection}{Factores limitantes}

Los factores limitantes son aquellos elementos o circunstancias que impiden el progreso y el logro de los objetivos establecidos en la empresa. Estos pueden ser internos, como la falta de recursos o capacidades del equipo, o externos, como la competencia intensa o cambios en el entorno económico. 

En el contexto del proyecto, estos factores representan barreras clave que deben gestionarse para garantizar la viabilidad del negocio

\begin{enumerate}
	\item \textbf{Costo de adquirir usuarios:} Como plataforma emergente en el ecosistema edtech bogotano, requeriremos una inversión significativa en marketing digital, incluyendo campañas en redes sociales. Por lo que los costos iniciales por usuario adquirido podrían ser elevados si las campañas no logran una alta efectividad.
	
	\item \textbf{Dependencia de oferta inicial:} La atractividad de la plataforma depende de una oferta inicial robusta de usuarios para fomentar interacciones comunitarias y generar datos para el machine learning. Sin un número suficiente de usuarios activos, será difícil crear el impulso viral necesario, lo que podría llevar a una adopción lenta y comprometer la sostenibilidad del modelo freemium.
	
	\item \textbf{Competencia intensa:} Competidores establecidos dominan el mercado con servicios de control financiero y banca digital. Existe el riesgo de que lancen módulos educativos similares, exigiendo una diferenciación continua y reforzar la propuesta de valor con el tiempo.
	
	\item \textbf{Brecha digital y demográfica:} Aunque Bogotá tiene alta penetración digital, persisten barreras en estratos bajos y trabajadores informales, con informalidad laboral alta, lo que podría limitar en cierto aspecto la adopción inicial y requiriendo estrategias de onboarding inclusivas.
\end{enumerate}